\section{A worked multi-step evidence dependency}\label{sec:C16:worked}
Consider a research agent asked whether a newly announced system improves long-horizon reliability. One source is a product announcement, another is an earlier technical report, and a third is a revised paper whose evaluation changes the reported gain. A whole-query embedding may rank the announcement first because it repeats the system name and the phrase ``long-horizon.'' The decisive evidence instead requires identifying the version lineage, the evaluation metric, and whether the tasks match the stated reliability claim.

A suitable evidence plan separates these needs. First resolve the canonical identity and versions. Next retrieve the method and evaluation protocol. Then locate the result table and limitations. Finally compare whether ``reliability'' means mean success, all-repeated-runs success, or a best-of-many metric. This is a constructed example, not a claim about a particular named system. It shows how retrieval targets can depend on earlier findings instead of being fully specified by the original query.

Represent the dependency as a directed acyclic graph of evidence requirements: identity precedes version-specific result selection; metric definition precedes interpretation; task distribution precedes transfer claims. The graph is a planning aid rather than a world model. A node can be unresolved, supported, contradicted, or unavailable. Reaching a final answer while an essential node remains unresolved should produce a qualified conclusion rather than a fabricated bridge.

Suppose the first retrieved result reports any-of-five success and the revised paper reports repeated-run reliability on a different split. Subtracting the two percentages would be invalid even if both were copied perfectly. Retrieval accuracy cannot repair a comparison-axis error. The synthesis stage must align attempts, denominators, splits, backbones, and budgets before ranking methods. A structured evidence table is often more reliable than a generated prose comparison because it makes missing fields explicit.

The proposed counterfactual test removes one requirement at a time. Hide the revision notice, hide the metric definition, or replace the primary report with a secondary summary. Observe which claim changes and whether the agent detects the missing dependency. A successful system should be able to revise its conclusion after new evidence arrives without discarding unrelated valid findings. This measures evidence integration and update behavior, not merely answer recall.

The same structure applies to execution. Before a tool action, the agent may need an object identifier, current resource version, relevant constraint, and authorization. These can come from different sources and have different freshness requirements. Grounded execution therefore extends retrieval beyond document question answering: it assembles a justified action under current state and authority, with the supporting evidence preserved in the run record.

